\documentclass[a11paper]{article}

\usepackage{karnaugh-map}
\usepackage{tabularx}
\usepackage{libs/titlepage}
\usepackage{amsmath}
\usepackage{amsthm}
\usepackage{libs/document}
\usepackage{booktabs}
\usepackage{multicol}
\usepackage{float}
\usepackage{csquotes}
\usepackage{fancyvrb}
\usepackage{hyperref}
\usepackage{tablefootnote}

\usepackage{witharrows}
\usepackage{enumitem}

\usepackage[usenames,dvipsnames]{xcolor}

\newcommand{\todo}[1]{\begin{color}{Red}\textbf{TODO:} #1\end{color}}
\newcommand{\note}[1]{\begin{color}{Orange}\textbf{NOTE:} #1\end{color}}
\newcommand{\fixme}[1]{\begin{color}{Fuchsia}\textbf{FIXME:} #1\end{color}}
\newcommand{\question}[1]{\begin{color}{ForestGreen}\textbf{QUESTION:} #1\end{color}}


% \renewcommand{\not}[1]{\overline{#1}}
\renewcommand{\not}[1]{#1^{\prime}}
\newcommand{\rhs}[1]{#1^{\prime}}
\newcommand{\xor}{\oplus}
\newcommand{\Cuz}[1]{\Arrow{\footnotesize #1}}

% \item $A+B=B+A \hspace{\fill}\text{commutativity}$
\usepackage{enumitem}       % customizable list environments

\newcommand{\cbox}[1][\relax]{\bgroup\setlength{\fboxsep}{4pt}\setlength{\fboxrule}{0.5pt}\fbox{}\egroup}

\title{Validation}

\class{Logique Séquentielle}
\classnb{APP1}

\teacher{Berthié Gouin-Ferland \& Nikola Zelovic}

\author{Benjamin Chausse --- CHAB1704}

\begin{document}
\maketitle

\section{Plan de tests}

\begin{table}[H]
  \centering
  \footnotesize
  \caption{Plan de vérification de la MEF de détection de fréquence}
  \begin{tabular}{llp{4.7cm}p{3cm}l}
    \toprule
    \multicolumn{3}{l}{Objectif ciblé} & \multicolumn{2}{l}{Valider la MEF}  \\
    & Condition à proscrire & Aucune &  \\
    \midrule
    \# & Test & Action & Résultats attendus & \\
    \midrule

    1 & Sinusoide pure surpassant les limites &
      Remplacer \Verb#input_signal.txt# par \Verb#sinusSignalHexa_26000Hz.txt# &
      La sortie alterne sans jamais converger & \cbox[name = hf_sin_over]{ } \\

    2 & Sinusoide pure de haute fréquence &
      Remplacer \Verb#input_signal.txt# par \Verb#sinusSignalHexa_20000Hz.txt# &
      La sortie converge aux alentours de deux\footnotemark[1] & \cbox[name = hf_sin_pure]{ } \\

    3 & Sinusoide pure de basse fréquence &
		  Remplacer \Verb#input_signal.txt# par \Verb#sinusSignalHexa_1000Hz.txt# &
      La sortie converge vers 48 après la première période\footnotemark[2] & \cbox[name = lf_sin_pure]{ } \\

    4 & Sinusoide obstruée de haute fréquence &
      Remplacer \Verb#input_signal.txt# par \Verb#sinusSignalHexa_20000Hz_Noise.txt# &
      La sortie converge vers deux\footnotemark[3] & \cbox[name = hf_sin_noise]{ } \\

    5 & Sinusoide obstruée de basse fréquence &
		  Remplacer \Verb#input_signal.txt# par \Verb#sinusSignalHexa_1000Hz_Noise.txt# &
      La sortie converge vers 48 après la première période\footnotemark[2]  & \cbox[name = lf_sin_noise]{ } \\

    6 & Onde Carré pure surpassant les limites &
      Remplacer \Verb#input_signal.txt# par \Verb#squareSignalHexa_26000Hz.txt# &
      Dû à la nature du carré, le lecteur converge vers une donnée erronée. & \cbox[name = hf_square_over]{ } \\

    7 & Onde Carrée pure de haute fréquence &
      Remplacer \Verb#input_signal.txt# par \Verb#squareSignalHexa_10000Hz.txt# &
      La sortie converge aux alentours de 4-5 après la première période\footnotemark[1] \\

    8 & Onde Carrée pure de basse fréquence &
		  Remplacer \Verb#input_signal.txt# par \Verb#squareSignalHexa_1000Hz.txt# &
      La sortie converge vers 48 après la première période\footnotemark[2]   & \cbox[name = lf_square_pure]{ } \\

    9 & Onde Carrée obstruée de haute fréquence &
      Remplacer \Verb#input_signal.txt# par \Verb#squareSignalHexa_10000Hz_Noise.txt# &
      La sortie converge aux alentours de 4-5\footnotemark[3] & \cbox[name = hf_square_noise]{ } \\

    10 & Onde Carrée obstruée de basse fréquence &
		  Remplacer \Verb#input_signal.txt# par \Verb#squareSignalHexa_1000Hz_Noise.txt# &
      La sortie converge vers 48 après la première  période (si le bruit est raisonable)&
      \cbox[name = lf_square_noise]{ } \\

    11 & Sinusoide de très basse fréquence &
		  Remplacer \Verb#input_signal.txt# par \Verb#sinusSignalHexa_200Hz.txt# &
      La sortie converge vers 240 après la première période\footnotemark[4]   & \cbox[name = lf_square_pure]{ } \\

    \bottomrule
  \end{tabular}
  \footnotetext[1]{hihihaha}
\end{table}

\footnotetext[1]{La sortie du module M5 est un nombre entier. Si le véritable
nombre d'échantillon par seconde est fractionnaire, la sortie entère va osciller
autour du nombre.}
\footnotetext[2]{Ce qui correspond à 30 en hexadécimal}
\footnotetext[3]{Le bruit impacte beaucoup les lectures à de haute fréquence.
On s'attends à une convergence générale mais pas parfaite dû à cette limitation}
\footnotetext[4]{Ce qui correspond à F0 en hexadécimal}




\end{document}
